\documentclass[11pt,a4paper]{article}
\usepackage[utf8]{inputenc}
\usepackage[T1]{fontenc}
\usepackage[spanish,es-minimal,es-noshorthands]{babel}
\AtBeginDocument{\renewcommand{\sin}{\operatorname{sen}}\renewcommand{\arcsin}{\operatorname{arcsen}}}
\usepackage{lmodern}
\usepackage[sfdefault,lining]{FiraSans}
\usepackage[a4paper,top=2.3cm,bottom=2.2cm,left=1.8cm,right=1.8cm,headheight=1.3cm,headsep=0.4cm,footskip=1cm]{geometry}
\usepackage{amsmath,amssymb}
\usepackage{icomma}
\usepackage{graphicx}
\usepackage[table]{xcolor}
\usepackage{tikz}
\usepackage{fontawesome5}
\usepackage{ccicons}
\usepackage{fancyhdr}
\usepackage{lastpage}
\usepackage{tasks}
\usepackage{needspace}
\usepackage[most]{tcolorbox}
\usepackage{array}
\usepackage[hidelinks]{hyperref}
\AtBeginDocument{\hypersetup{pdftitle={Ficha de ejercicios de la unidad 5: Trigonometría. 1º Bach.},pdfsubject={Matemáticas I 1º Bach.},pdfkeywords={trigonometría, radianes, razones trigonométricas, circunferencia goniométrica, resolución de triángulos, teorema del seno, teorema del coseno, fórmulas trigonométricas, ecuaciones trigonométricas, funciones trigonométricas}}}

\definecolor{navy}{HTML}{1B2748}
\definecolor{accent}{HTML}{E8603A}
\definecolor{teal}{HTML}{1F9E90}
\definecolor{soft}{HTML}{EEF1F8}
\definecolor{muted}{HTML}{5D6685}

\setlength{\parindent}{0pt}
\setlength{\parskip}{3pt}
\settasks{label=\alph*), label-format=\bfseries\color{navy}, label-width=1.3em, item-indent=1.8em, column-sep=0.6em, after-item-skip=5pt, before-skip=3pt, after-skip=4pt}

% ---------- capçalera i peu ----------
\newcommand{\llicencia}{\href{https://creativecommons.org/licenses/by-nc-sa/4.0/deed.es}{\ccbyncsa}\ \ Licencia Creative Commons Atribución-NoComercial-CompartirIgual 4.0 Internacional (CC BY-NC-SA 4.0)}
\pagestyle{fancy}
\fancyhf{}
\fancyhead[L]{\raisebox{-0.35\height}{\includegraphics[height=0.95cm]{logo-gva.png}}}
\fancyhead[C]{\small\color{muted}Ficha de ejercicios · Unidad 5: Trigonometría · 1\textsuperscript{o} Bach.}
\fancyhead[R]{\raisebox{-0.35\height}{\includegraphics[height=0.95cm]{logo-vermellar.png}}}
\fancyfoot[L]{\scriptsize\color{muted}\llicencia}
\fancyfoot[R]{\small\color{muted}\thepage}
\renewcommand{\headrulewidth}{0.4pt}
\renewcommand{\headrule}{\color{muted!40}\hrule height \headrulewidth}
\fancypagestyle{primera}{\fancyhf{}\renewcommand{\headrulewidth}{0pt}
  \fancyfoot[L]{\scriptsize\color{muted}\llicencia}
  \fancyfoot[R]{\small\color{muted}\thepage}}

% ---------- nivells i exercicis ----------
\newcommand{\estrelles}[1]{{\color{accent}\ifcase#1\or\faStar\faStar[regular]\faStar[regular]\or\faStar\faStar\faStar[regular]\or\faStar\faStar\faStar\fi}}
\newcounter{exer}
\newcommand{\ex}[1]{\par\needspace{4\baselineskip}\medskip\refstepcounter{exer}\noindent{\color{navy}\bfseries\theexer.}\ \ }

\newcommand{\seccio}[2]{\par\needspace{16\baselineskip}\bigskip
  \begin{tcolorbox}[colback=navy,colframe=navy,coltext=white,arc=2mm,boxrule=0pt,left=3mm,right=3mm,top=1.2mm,bottom=1.2mm]
  {\large\bfseries #1\quad #2}
  \end{tcolorbox}\vspace{-2pt}}

\newcommand{\exf}[1]{\par\medskip\refstepcounter{exer}\noindent{\color{navy}\bfseries\theexer.}\ \ }
\usetikzlibrary{calc}
\newcolumntype{C}[1]{>{\centering\arraybackslash}p{#1}}
% quadre de recordatori (només a la fitxa d'exercicis)
\newcommand{\recorda}[1]{\par\smallskip\begin{tcolorbox}[colback=accent!7,colframe=accent!55,arc=2mm,boxrule=0.5pt,left=2.5mm,right=2.5mm,top=1mm,bottom=1mm]\small{\color{accent}\faLightbulb[regular]}\ \textbf{Recuerda:}\ #1\end{tcolorbox}}
\newcommand{\res}[1]{#1}\newcommand{\rest}[1]{#1}

\def\fmneed{11}% que el títol d'una secció no quede sol al final del full
% =====================================================================
% Capçalera FentMates de les fitxes: un color per a cada curs.
% Els colors són els de la pàgina inicial (LEVELS a assets/js/materials.js);
% el logo ve de _eines/logo/pdf/ (python3 _eines/logo/genera.py).
%
% Ús, al final del preàmbul de la unitat:
%   % =====================================================================
% Capçalera FentMates de les fitxes: un color per a cada curs.
% Els colors són els de la pàgina inicial (LEVELS a assets/js/materials.js);
% el logo ve de _eines/logo/pdf/ (python3 _eines/logo/genera.py).
%
% Ús, al final del preàmbul de la unitat:
%   % =====================================================================
% Capçalera FentMates de les fitxes: un color per a cada curs.
% Els colors són els de la pàgina inicial (LEVELS a assets/js/materials.js);
% el logo ve de _eines/logo/pdf/ (python3 _eines/logo/genera.py).
%
% Ús, al final del preàmbul de la unitat:
%   \input{../comu/fentmates.tex}
%   \fmcurs{1eso}
% i, al començament del document:
%   \capcalera{Matemàtiques}{1r ESO}{Fitxa d'exercicis · Unitat 2}{Divisibilitat}
% Opcional abans de \input: \def\fmneed{14} (espai mínim abans de cada secció).
% =====================================================================
\makeatletter
\def\fm@logodir{../../logo/pdf}
\providecommand{\fmneed}{10}
\def\fm@cursnom{}

% clau del curs, color i color del text damunt del color
\newcommand{\fm@curs}[3]{\@namedef{fm@color@#1}{#2}\@namedef{fm@text@#1}{#3}}
\fm@curs{1eso}{E6B84A}{navy}
\fm@curs{2eso}{3B82F6}{white}
\fm@curs{3eso}{FF7A50}{white}
\fm@curs{4eso-a}{A78BFA}{white}
\fm@curs{4eso-b}{6D28D9}{white}
\fm@curs{1batx-mi}{9B1C1C}{white}
\fm@curs{1batx-ccss}{EF6B6B}{white}
\fm@curs{2batx-mii}{1E6B47}{white}
\fm@curs{2batx-ccss}{2BB3A3}{white}

\newcommand{\fmcurs}[1]{%
  \@ifundefined{fm@color@#1}{\PackageError{fentmates}{Curs desconegut: #1}{}}{}%
  \edef\fm@tmp{\noexpand\definecolor{curs}{HTML}{\@nameuse{fm@color@#1}}%
               \noexpand\colorlet{curstext}{\@nameuse{fm@text@#1}}}\fm@tmp
  \colorlet{cursfosc}{curs!72!black}%
  \def\fm@logo{\fm@logodir/vertical-#1.pdf}%
  \def\fm@marca{\fm@logodir/marca-#1.pdf}}

% ---------- capçalera de la primera pàgina ----------
% #1 matèria (petit), #2 curs, #3 tipus de document i unitat, #4 títol
\newcommand{\capcalera}[4]{%
  \gdef\fm@cursnom{#2}%
  \par\noindent
  \begin{tcolorbox}[colback=curs,colframe=curs,boxrule=0pt,arc=2.2mm,width=3.5cm,height=2.15cm,
      valign=center,left=3mm,right=2mm,top=1mm,bottom=1mm,nobeforeafter,box align=center]
    \color{curstext}{\footnotesize #1}\par\vspace{1pt}{\fontsize{19}{21}\selectfont\bfseries #2\par}%
  \end{tcolorbox}\hspace{5mm}%
  \begin{minipage}[c]{\dimexpr\linewidth-3.5cm-5mm-2.9cm\relax}
    {\color{cursfosc}\small\bfseries\MakeUppercase{#3}}\par\vspace{2pt}
    {\color{navy}\fontsize{20}{23}\selectfont\bfseries #4\par}
  \end{minipage}\hfill
  \begin{minipage}[c]{2.6cm}\raggedleft\includegraphics[width=2.6cm]{\fm@logo}\end{minipage}\par
  \vspace{2pt}}

% ---------- barres de secció del color del curs ----------
\renewcommand{\seccio}[2]{\par\needspace{\fmneed\baselineskip}\bigskip
  \begin{tcolorbox}[colback=curs,colframe=curs,coltext=curstext,arc=2mm,boxrule=0pt,left=3mm,right=3mm,top=1.2mm,bottom=1.2mm]
  {\large\bfseries #1\quad #2}
  \end{tcolorbox}\vspace{-2pt}}

% ---------- estrelles de dificultat del color del curs (a Batxillerat no se'n posen) ----------
\@ifundefined{estrelles}{}{%
  \renewcommand{\estrelles}[1]{{\color{curs}\ifcase#1\or\faStar\faStar[regular]\faStar[regular]\or\faStar\faStar\faStar[regular]\or\faStar\faStar\faStar\fi}}}

% ---------- capçalera i peu de les pàgines ----------
% primera pàgina: logos institucionals al peu
\fancypagestyle{primera}{\fancyhf{}\renewcommand{\headrulewidth}{0pt}
  \fancyfoot[L]{\raisebox{-0.32\height}{\includegraphics[height=0.78cm]{logo-gva.png}}\hspace{5mm}%
    \raisebox{-0.32\height}{\includegraphics[height=0.78cm]{logo-vermellar.png}}}
  \fancyfoot[R]{\scriptsize\color{muted}\href{https://creativecommons.org/licenses/by-nc-sa/4.0/deed.ca}{\ccbyncsa}\ \ CC BY-NC-SA 4.0 · fentmates.org\qquad\small\thepage}}
% la resta: el símbol a l'esquerra i el curs a la dreta, amb el títol al centre com abans
\fancyhead[L]{\raisebox{-0.3\height}{\includegraphics[height=0.62cm]{\fm@marca}}}
\fancyhead[R]{\small\bfseries\color{cursfosc}\fm@cursnom}
\renewcommand{\headrule}{\color{curs!45}\hrule height \headrulewidth}
\makeatother

%   \fmcurs{1eso}
% i, al començament del document:
%   \capcalera{Matemàtiques}{1r ESO}{Fitxa d'exercicis · Unitat 2}{Divisibilitat}
% Opcional abans de \input: \def\fmneed{14} (espai mínim abans de cada secció).
% =====================================================================
\makeatletter
\def\fm@logodir{../../logo/pdf}
\providecommand{\fmneed}{10}
\def\fm@cursnom{}

% clau del curs, color i color del text damunt del color
\newcommand{\fm@curs}[3]{\@namedef{fm@color@#1}{#2}\@namedef{fm@text@#1}{#3}}
\fm@curs{1eso}{E6B84A}{navy}
\fm@curs{2eso}{3B82F6}{white}
\fm@curs{3eso}{FF7A50}{white}
\fm@curs{4eso-a}{A78BFA}{white}
\fm@curs{4eso-b}{6D28D9}{white}
\fm@curs{1batx-mi}{9B1C1C}{white}
\fm@curs{1batx-ccss}{EF6B6B}{white}
\fm@curs{2batx-mii}{1E6B47}{white}
\fm@curs{2batx-ccss}{2BB3A3}{white}

\newcommand{\fmcurs}[1]{%
  \@ifundefined{fm@color@#1}{\PackageError{fentmates}{Curs desconegut: #1}{}}{}%
  \edef\fm@tmp{\noexpand\definecolor{curs}{HTML}{\@nameuse{fm@color@#1}}%
               \noexpand\colorlet{curstext}{\@nameuse{fm@text@#1}}}\fm@tmp
  \colorlet{cursfosc}{curs!72!black}%
  \def\fm@logo{\fm@logodir/vertical-#1.pdf}%
  \def\fm@marca{\fm@logodir/marca-#1.pdf}}

% ---------- capçalera de la primera pàgina ----------
% #1 matèria (petit), #2 curs, #3 tipus de document i unitat, #4 títol
\newcommand{\capcalera}[4]{%
  \gdef\fm@cursnom{#2}%
  \par\noindent
  \begin{tcolorbox}[colback=curs,colframe=curs,boxrule=0pt,arc=2.2mm,width=3.5cm,height=2.15cm,
      valign=center,left=3mm,right=2mm,top=1mm,bottom=1mm,nobeforeafter,box align=center]
    \color{curstext}{\footnotesize #1}\par\vspace{1pt}{\fontsize{19}{21}\selectfont\bfseries #2\par}%
  \end{tcolorbox}\hspace{5mm}%
  \begin{minipage}[c]{\dimexpr\linewidth-3.5cm-5mm-2.9cm\relax}
    {\color{cursfosc}\small\bfseries\MakeUppercase{#3}}\par\vspace{2pt}
    {\color{navy}\fontsize{20}{23}\selectfont\bfseries #4\par}
  \end{minipage}\hfill
  \begin{minipage}[c]{2.6cm}\raggedleft\includegraphics[width=2.6cm]{\fm@logo}\end{minipage}\par
  \vspace{2pt}}

% ---------- barres de secció del color del curs ----------
\renewcommand{\seccio}[2]{\par\needspace{\fmneed\baselineskip}\bigskip
  \begin{tcolorbox}[colback=curs,colframe=curs,coltext=curstext,arc=2mm,boxrule=0pt,left=3mm,right=3mm,top=1.2mm,bottom=1.2mm]
  {\large\bfseries #1\quad #2}
  \end{tcolorbox}\vspace{-2pt}}

% ---------- estrelles de dificultat del color del curs (a Batxillerat no se'n posen) ----------
\@ifundefined{estrelles}{}{%
  \renewcommand{\estrelles}[1]{{\color{curs}\ifcase#1\or\faStar\faStar[regular]\faStar[regular]\or\faStar\faStar\faStar[regular]\or\faStar\faStar\faStar\fi}}}

% ---------- capçalera i peu de les pàgines ----------
% primera pàgina: logos institucionals al peu
\fancypagestyle{primera}{\fancyhf{}\renewcommand{\headrulewidth}{0pt}
  \fancyfoot[L]{\raisebox{-0.32\height}{\includegraphics[height=0.78cm]{logo-gva.png}}\hspace{5mm}%
    \raisebox{-0.32\height}{\includegraphics[height=0.78cm]{logo-vermellar.png}}}
  \fancyfoot[R]{\scriptsize\color{muted}\href{https://creativecommons.org/licenses/by-nc-sa/4.0/deed.ca}{\ccbyncsa}\ \ CC BY-NC-SA 4.0 · fentmates.org\qquad\small\thepage}}
% la resta: el símbol a l'esquerra i el curs a la dreta, amb el títol al centre com abans
\fancyhead[L]{\raisebox{-0.3\height}{\includegraphics[height=0.62cm]{\fm@marca}}}
\fancyhead[R]{\small\bfseries\color{cursfosc}\fm@cursnom}
\renewcommand{\headrule}{\color{curs!45}\hrule height \headrulewidth}
\makeatother

%   \fmcurs{1eso}
% i, al començament del document:
%   \capcalera{Matemàtiques}{1r ESO}{Fitxa d'exercicis · Unitat 2}{Divisibilitat}
% Opcional abans de \input: \def\fmneed{14} (espai mínim abans de cada secció).
% =====================================================================
\makeatletter
\def\fm@logodir{../../logo/pdf}
\providecommand{\fmneed}{10}
\def\fm@cursnom{}

% clau del curs, color i color del text damunt del color
\newcommand{\fm@curs}[3]{\@namedef{fm@color@#1}{#2}\@namedef{fm@text@#1}{#3}}
\fm@curs{1eso}{E6B84A}{navy}
\fm@curs{2eso}{3B82F6}{white}
\fm@curs{3eso}{FF7A50}{white}
\fm@curs{4eso-a}{A78BFA}{white}
\fm@curs{4eso-b}{6D28D9}{white}
\fm@curs{1batx-mi}{9B1C1C}{white}
\fm@curs{1batx-ccss}{EF6B6B}{white}
\fm@curs{2batx-mii}{1E6B47}{white}
\fm@curs{2batx-ccss}{2BB3A3}{white}

\newcommand{\fmcurs}[1]{%
  \@ifundefined{fm@color@#1}{\PackageError{fentmates}{Curs desconegut: #1}{}}{}%
  \edef\fm@tmp{\noexpand\definecolor{curs}{HTML}{\@nameuse{fm@color@#1}}%
               \noexpand\colorlet{curstext}{\@nameuse{fm@text@#1}}}\fm@tmp
  \colorlet{cursfosc}{curs!72!black}%
  \def\fm@logo{\fm@logodir/vertical-#1.pdf}%
  \def\fm@marca{\fm@logodir/marca-#1.pdf}}

% ---------- capçalera de la primera pàgina ----------
% #1 matèria (petit), #2 curs, #3 tipus de document i unitat, #4 títol
\newcommand{\capcalera}[4]{%
  \gdef\fm@cursnom{#2}%
  \par\noindent
  \begin{tcolorbox}[colback=curs,colframe=curs,boxrule=0pt,arc=2.2mm,width=3.5cm,height=2.15cm,
      valign=center,left=3mm,right=2mm,top=1mm,bottom=1mm,nobeforeafter,box align=center]
    \color{curstext}{\footnotesize #1}\par\vspace{1pt}{\fontsize{19}{21}\selectfont\bfseries #2\par}%
  \end{tcolorbox}\hspace{5mm}%
  \begin{minipage}[c]{\dimexpr\linewidth-3.5cm-5mm-2.9cm\relax}
    {\color{cursfosc}\small\bfseries\MakeUppercase{#3}}\par\vspace{2pt}
    {\color{navy}\fontsize{20}{23}\selectfont\bfseries #4\par}
  \end{minipage}\hfill
  \begin{minipage}[c]{2.6cm}\raggedleft\includegraphics[width=2.6cm]{\fm@logo}\end{minipage}\par
  \vspace{2pt}}

% ---------- barres de secció del color del curs ----------
\renewcommand{\seccio}[2]{\par\needspace{\fmneed\baselineskip}\bigskip
  \begin{tcolorbox}[colback=curs,colframe=curs,coltext=curstext,arc=2mm,boxrule=0pt,left=3mm,right=3mm,top=1.2mm,bottom=1.2mm]
  {\large\bfseries #1\quad #2}
  \end{tcolorbox}\vspace{-2pt}}

% ---------- estrelles de dificultat del color del curs (a Batxillerat no se'n posen) ----------
\@ifundefined{estrelles}{}{%
  \renewcommand{\estrelles}[1]{{\color{curs}\ifcase#1\or\faStar\faStar[regular]\faStar[regular]\or\faStar\faStar\faStar[regular]\or\faStar\faStar\faStar\fi}}}

% ---------- capçalera i peu de les pàgines ----------
% primera pàgina: logos institucionals al peu
\fancypagestyle{primera}{\fancyhf{}\renewcommand{\headrulewidth}{0pt}
  \fancyfoot[L]{\raisebox{-0.32\height}{\includegraphics[height=0.78cm]{logo-gva.png}}\hspace{5mm}%
    \raisebox{-0.32\height}{\includegraphics[height=0.78cm]{logo-vermellar.png}}}
  \fancyfoot[R]{\scriptsize\color{muted}\href{https://creativecommons.org/licenses/by-nc-sa/4.0/deed.ca}{\ccbyncsa}\ \ CC BY-NC-SA 4.0 · fentmates.org\qquad\small\thepage}}
% la resta: el símbol a l'esquerra i el curs a la dreta, amb el títol al centre com abans
\fancyhead[L]{\raisebox{-0.3\height}{\includegraphics[height=0.62cm]{\fm@marca}}}
\fancyhead[R]{\small\bfseries\color{cursfosc}\fm@cursnom}
\renewcommand{\headrule}{\color{curs!45}\hrule height \headrulewidth}
\makeatother

\fmcurs{1batx-mi}

\begin{document}
\thispagestyle{primera}

\capcalera{Matemáticas I}{1\textsuperscript{o} Bach.}{Ficha de ejercicios · Unidad 5}{\raggedright Trigonometría}
\vspace{10pt}
Nombre y apellidos: \rule{7.8cm}{0.4pt}\hfill Grupo: \rule{1.6cm}{0.4pt}\hfill Fecha: \rule{2.2cm}{0.4pt}
\vspace{8pt}

\begin{tcolorbox}[colback=soft,colframe=soft,arc=2mm,left=3mm,right=3mm,top=1.5mm,bottom=1.5mm]
\small Resuelve los ejercicios en el cuaderno mostrando todos los pasos y haz un dibujo siempre que te ayude. Comprueba el modo de la calculadora (grados o radianes) y, en las ecuaciones, comprueba siempre las soluciones. Los ejercicios 1--41 son los de la ficha «Trigonometría 1» del departamento y los 42--60, los de «Trigonometría 2», en el mismo orden; la sección 13 tiene otros nuevos para ampliar y aplicar lo que has aprendido.
\end{tcolorbox}

% =====================================================
\seccio{1}{Medida de ángulos}

\exf{1} Transforma estos ángulos en grados o radianes, según corresponda.\begin{tasks}(4)
  \task $225^\circ$
  \task $75^\circ$
  \task $160^\circ$
  \task $-1{,}5$ rad
  \task $2$ rad
  \task $540^\circ$
  \task $-270^\circ$
  \task $140^\circ40'$
  \task $\frac{\pi}{3}$ rad
  \task $0{,}3$ rad
  \task $-\frac{6\pi}{5}$ rad
  \task $315^\circ$
\end{tasks}


% =====================================================
\seccio{2}{Trigonometría con la calculadora}

\exf{1} Con la calculadora, halla estas razones (con cuatro decimales). Comprueba antes si debe estar en modo \texttt{D} (grados) o \texttt{R} (radianes).\begin{tasks}(3)
  \task $\sin(319^\circ12'52'')$
  \task $\cos(434^\circ26')$
  \task $\text{tg}\,\frac{11\pi}{8}$
  \task $\sec135^\circ$
  \task $\text{tg}\,7{,}03$
  \task $\text{cosec}\,5{,}24$
  \task $\sin\frac{8\pi}{5}$
  \task $\text{cosec}\,(55'45'')$
  \task $\cos\frac{\pi}{3}$
  \task $\text{cosec}\,(200^\circ16')$
  \task $\sec\frac{5\pi}{4}$
  \task $\text{tg}\,(183^\circ13'53'')$
  \task $\sin(333^\circ0'55'')$
  \task $\cos3{,}845$
  \task $\text{cotg}\,\frac{11\pi}{6}$
\end{tasks}

\ex{1} Halla los ángulos de $[0^\circ,360^\circ)$ que cumplen:\begin{tasks}(4)
  \task $\cos x=0{,}4539$
  \task $\sin x=0{,}9284$
  \task $\text{tg}\,x=-0{,}5459$
  \task $\cos x=-0{,}3996$
  \task $\text{cotg}\,x=2{,}1618$
  \task $\cos x=-0{,}2926$
  \task $\sin x=0{,}6862$
  \task $\sin x=-0{,}3308$
\end{tasks}\recorda{la calculadora da un solo ángulo. El otro: $180^\circ-\alpha$ (seno), $360^\circ-\alpha$ (coseno) o $\alpha+180^\circ$ (tangente).}


% =====================================================
\seccio{3}{Razones trigonométricas de un ángulo agudo}

\exf{1} Completa la tabla siguiente sin calculadora:\begin{center}\renewcommand{\arraystretch}{1.5}\begin{tabular}{|l|*{5}{C{1.55cm}|}}\hline
\textbf{Grados} & $0^\circ$ & $30^\circ$ & $45^\circ$ & $60^\circ$ & $90^\circ$\\\hline
\textbf{Radianes} &  &  &  &  & \\\hline
$\sin$ &  &  &  &  & \\\hline $\cos$ &  &  &  &  & \\\hline $\text{tg}$ &  &  &  &  & \\\hline $\text{cosec}$ &  &  &  &  & \\\hline $\sec$ &  &  &  &  & \\\hline $\text{cotg}$ &  &  &  &  & \\\hline
\end{tabular}\end{center}

\ex{1} Utiliza las relaciones fundamentales para hallar las otras razones trigonométricas de los ángulos agudos $\alpha$, $\beta$ y $\gamma$:\begin{tasks}(3)
  \task $\cos\alpha=\frac{\sqrt5}{3}$
  \task $\sin\beta=\frac35$
  \task $\text{tg}\,\gamma=3$
  \task $\cos\alpha=0{,}86$
  \task $\text{tg}\,\beta=2{,}83$
  \task $\text{tg}\,\gamma=1$
\end{tasks}\recorda{$\sin^2\alpha+\cos^2\alpha=1$, \ $\text{tg}\,\alpha=\dfrac{\sin\alpha}{\cos\alpha}$, \ $1+\text{tg}^2\alpha=\dfrac{1}{\cos^2\alpha}$.}


% =====================================================
\seccio{4}{Razones trigonométricas de ángulos cualesquiera}

\exf{1} Completa la tabla siguiente sin calculadora:\begin{center}\renewcommand{\arraystretch}{1.45}\begin{tabular}{|C{1.4cm}|C{1.6cm}|C{1.6cm}|C{1.6cm}|C{1.6cm}|}\hline\textbf{Grados} & \textbf{Radianes} & $\sin$ & $\cos$ & $\text{tg}$\\\hline$120^\circ$ & & & & \\\hline $135^\circ$ & & & & \\\hline $150^\circ$ & & & & \\\hline $180^\circ$ & & & & \\\hline $210^\circ$ & & & & \\\hline $225^\circ$ & & & & \\\hline $240^\circ$ & & & & \\\hline $270^\circ$ & & & & \\\hline $300^\circ$ & & & & \\\hline $315^\circ$ & & & & \\\hline $330^\circ$ & & & & \\\hline $360^\circ$ & & & & \\\hline \end{tabular}\end{center}

\ex{1} Sin utilizar la calculadora, calcula:\begin{tasks}(2)
  \task $\sin\pi-\text{tg}\,\pi+2\cos\pi$
  \task $\cos\frac{\pi}{2}-\sin\frac{3\pi}{2}$
  \task $\text{tg}\,2\pi-\sin\frac{\pi}{2}-2\cos\frac{3\pi}{2}$
  \task $\cos2\pi-\sin2\pi$
\end{tasks}

\ex{2} La tabla muestra razones trigonométricas de ángulos de cuadrantes diferentes. Sin determinarlos, complétala con las razones que faltan.\begin{center}\renewcommand{\arraystretch}{1.3}\begin{tabular}{|l|C{2.4cm}|C{2.4cm}|C{2.4cm}|}\hline\textbf{Cuadrante} & $\sin$ & $\cos$ & $\text{tg}$\\\hline Segundo & $0{,}6702$ &  & \\\hline Tercero &  & $-0{,}4539$ & \\\hline Cuarto &  &  & $-0{,}7459$\\\hline Tercero & $-0{,}7822$ &  & \\\hline Segundo &  &  & $-1{,}9004$\\\hline Cuarto &  & $0{,}6983$ & \\\hline \end{tabular}\end{center}

\ex{1} Calcula el seno, el coseno y la tangente del ángulo $\alpha$:\begin{tasks}(2)
  \task $\sec\alpha=4$ \ y \ $270^\circ<\alpha<360^\circ$
  \task $\text{cotg}\,\alpha=-\frac12$ \ y \ $\frac{\pi}{2}<\alpha<\pi$
\end{tasks}

\ex{1} Sabiendo que el ángulo $\alpha$ está en el segundo cuadrante ($90^\circ<\alpha<180^\circ$) y $\sin\alpha=0{,}62$, calcula las otras razones trigonométricas.

\ex{1} Sabiendo que el ángulo $\alpha$ está en el tercer cuadrante ($180^\circ<\alpha<270^\circ$) y $\cos\alpha=-0{,}83$, calcula las otras razones trigonométricas.

\ex{1} Sabiendo que el ángulo $\alpha$ está en el cuarto cuadrante ($270^\circ<\alpha<360^\circ$) y $\text{tg}\,\alpha=-0{,}92$, calcula las otras razones trigonométricas.

\ex{1} Determina el valor exacto de estas razones trigonométricas:\begin{tasks}(4)
  \task $\sin13290^\circ$
  \task $\cos(-1680^\circ)$
  \task $\text{tg}\,3825^\circ$
  \task $\cos4995^\circ$
  \task $\sin(-1710^\circ)$
  \task $\text{tg}\,3630^\circ$
  \task $\cos(-36000^\circ)$
  \task $\sin(-330^\circ)$
\end{tasks}


% =====================================================
\seccio{5}{Relaciones entre las razones trigonométricas de algunos ángulos}

\exf{2} \begin{tasks}(1)\task Halla $\sin\alpha$, $\cos\alpha$, $\sin(180^\circ-\alpha)$, $\cos(90^\circ+\alpha)$ y $\text{tg}\,(180^\circ+\alpha)$, sin calcular el ángulo $\alpha$, sabiendo que $\text{tg}\,\alpha=-\frac{\sqrt5}{2}$ y $90^\circ<\alpha<180^\circ$. \task Halla con la calculadora el valor del ángulo $\alpha$.\end{tasks}

\ex{1} Sabiendo que $\cos\alpha=0{,}8$ y $\sin\alpha=0{,}6$, calcula:\begin{tasks}(3)
  \task $\cos(180^\circ+\alpha)$
  \task $\sin(180^\circ-\alpha)$
  \task $\text{tg}\,(-\alpha)$
  \task $\sin(90^\circ-\alpha)$
  \task $\cos(90^\circ+\alpha)$
  \task $\text{tg}\,(90^\circ-\alpha)$
\end{tasks}

\ex{1} Si $\sin42^\circ=0{,}67$, $\cos42^\circ=0{,}74$ y $\text{tg}\,42^\circ=0{,}9$, halla las razones trigonométricas siguientes sin utilizar la calculadora:\begin{tasks}(3)
  \task $\cos48^\circ$
  \task $\sin(-48^\circ)$
  \task $\sin138^\circ$
  \task $\text{tg}\,318^\circ$
  \task $\cos222^\circ$
  \task $\text{tg}\,858^\circ$
\end{tasks}


% =====================================================
\seccio{6}{Resolución de triángulos rectángulos}

\exf{1} Los triángulos siguientes son rectángulos y se llaman $ABC$, con el ángulo recto en $C$. Resuélvelos:\begin{tasks}(1)
  \task Datos: $c=32$ cm, $\hat B=57^\circ$. Calcula $a$ y $b$.
  \task Datos: $a=250$ m, $b=308$ m. Calcula $c$ y $\hat A$.
  \task Datos: $a=35$ cm, $\hat A=32^\circ$. Calcula $b$ y $c$.
\end{tasks}

\ex{1} Un mástil está sujeto al suelo con dos cables de 12 m que forman ángulos de $50^\circ$ con el suelo. Calcula la altura del mástil y la distancia de la base a los puntos de sujeción.

\ex{2} El tejado de una casa tiene la forma y las medidas que se indican en la figura. Calcula la distancia $\overline{AB}$.\begin{center}\begin{tikzpicture}[scale=0.62]
\coordinate (A) at (0,0); \coordinate (P) at ($(A)+(60:4)$); \coordinate (Q) at ($(P)+(-30:1.7)$); \coordinate (R) at ($(Q)+(20:3.5)$); \coordinate (B) at ($(R|-A)+(2.7,0)$);
\draw[muted, dashed] (A)--(B); \draw[muted, dashed] (P)--++(1.6,0); \draw[muted, dashed] (Q)--++(1.5,0);
\draw[navy, line width=2.2pt, line join=round] (A)--(P)--(Q)--(R)--(B);
\node[below left] at (A) {$A$}; \node[below right] at (B) {$B$};
\node[font=\small, rotate=60] at ($(A)!0.5!(P)+(-0.4,0.22)$) {4 m}; \node[font=\small, rotate=-30] at ($(P)!0.5!(Q)+(-0.25,-0.45)$) {1,7 m};
\node[font=\small, rotate=20] at ($(Q)!0.5!(R)+(-0.15,0.42)$) {3,5 m}; \node[font=\small, rotate=-55] at ($(R)!0.5!(B)+(0.45,0.3)$) {4,2 m};
\draw (A)++(0.7,0) arc (0:60:0.7); \node[font=\scriptsize] at (1.05,0.4) {$60^\circ$};
\draw (P)++(0.75,0) arc (0:-30:0.75); \node[font=\scriptsize] at ($(P)+(1.35,-0.2)$) {$30^\circ$};
\draw (Q)++(0.75,0) arc (0:20:0.75); \node[font=\scriptsize] at ($(Q)+(1.3,0.12)$) {$20^\circ$};
\draw (B)++(-0.7,0) arc (180:125:0.7); \node[font=\scriptsize] at ($(B)+(-1.05,0.35)$) {$50^\circ$};
\end{tikzpicture}\end{center}

\ex{2} Un avión despega formando un ángulo de $15^\circ$ con el suelo. A 6 km del punto de salida (en horizontal) aumenta $10^\circ$ la inclinación, de manera que sube con un ángulo de $25^\circ$ con el suelo. Una montaña de 2\,500 m se encuentra a 11 km del punto de partida. ¿Volará el avión por encima de la montaña o tendrá que cambiar la dirección?

\ex{2} En un triángulo isósceles, el lado desigual mide 10 m y el ángulo opuesto es de $40^\circ$. Calcula su perímetro y su área.

\ex{1} Calcula el área de un paralelogramo cuyos lados, de 9 cm y 16 cm, forman un ángulo de $65^\circ$.


% =====================================================
\seccio{7}{Teorema del seno y teorema del coseno}

\exf{2} Resuelve los triángulos siguientes:\begin{tasks}(1)
  \task $a=10$ cm, \ $b=14$ cm, \ $c=8$ cm
  \task $b=6$ cm, \ $c=9$ cm, \ $\hat A=39^\circ12'$
  \task $a=7$ cm, \ $\hat B=38^\circ49'$, \ $\hat C=66^\circ40'$
  \task $a=8$ cm, \ $b=10$ cm, \ $\hat B=36^\circ38'$
  \task $a=9$ cm, \ $c=5$ cm, \ $\hat B=103^\circ27'$
\end{tasks}

\ex{1} Decide si las medidas siguientes corresponden a las longitudes de los lados de un triángulo y, si es así, indica si es acutángulo, rectángulo u obtusángulo.\begin{tasks}(4)
  \task 12, 11 y 9 cm
  \task 23, 14 y 8 cm
  \task 26, 24 y 10 cm
  \task 40, 30 y 20 m
\end{tasks}


% =====================================================
\seccio{8}{Resolución de triángulos: problemas}

\exf{1} Una barca está atada a la orilla de un canal con una cuerda de 8 metros. En un determinado momento, esta cuerda forma un ángulo de $38^\circ$ con la orilla. ¿A qué distancia de la orilla está la barca?

\ex{1} Si nos situamos a 40 metros de la chimenea de una fábrica, la vemos bajo un ángulo de $26^\circ$. ¿Qué altura tiene? Considera que los ojos del observador están situados a 175 cm del suelo.

\ex{2} La anchura de un escenario de teatro es de 8 metros. Las localidades que hemos comprado están situadas a una distancia de 6 m y de 12 m de cada uno de los extremos laterales del escenario. ¿Cuál es el ángulo de visión que tendremos para ver la representación?

\ex{2} Desde un punto del suelo medimos el ángulo bajo el que se ve un edificio y obtenemos $40^\circ$. Nos alejamos 30 m y el ángulo es ahora de $28^\circ$. Calcula la altura del edificio y la distancia desde la que se hizo la primera observación.

\ex{3} Desde una torre de vigilancia de 25 m de altura observamos dos árboles situados en orillas opuestas de un río bajo un ángulo de $15^\circ$. Los dos árboles están alineados con el pie de la torre y la distancia de esta al río es de 50 m. Calcula la anchura del río.

\ex{1} El radar de un barco detecta un objeto no identificado a 40 m de profundidad y en una dirección que forma $15^\circ$ con la horizontal. ¿Qué distancia tiene que recorrer un buzo para llegar desde el barco hasta el objeto?

\ex{2} Un árbol se sujeta al suelo con dos cables de acero, uno a cada lado, que forman ángulos de $43^\circ$ y de $57^\circ50'$ con el suelo. Si las distancias de los cables al pie del árbol suman 15 m, ¿cuál es la altura del árbol?

\ex{2} El área de un triángulo rectángulo es de 28 cm$^2$ y uno de sus ángulos mide $60^\circ$.\begin{tasks}(1)
  \task ¿Cuánto mide cada uno de sus ángulos?
  \task Calcula la longitud de los lados y el perímetro del triángulo.
\end{tasks}

\ex{1} Salvador mide 1,70 m y observa que su sombra es de 50 cm a una determinada hora del día. ¿Con qué inclinación llegan los rayos solares a esa hora?

\ex{2} Una casa de planta rectangular mide 12 m de largo y 8 m de ancho. El techo, con una inclinación de $18^\circ$, es una superficie plana inclinada cuya parte más elevada está situada sobre uno de los lados mayores del rectángulo. Calcula el área del tejado.

\ex{2} Dos exploradores se han perdido y deciden seguir caminos diferentes para conseguir ayuda; sus trayectorias forman un ángulo de $54^\circ$. Si uno camina a 5 km/h y el otro a 4 km/h, ¿a qué distancia están al cabo de 2 horas? ¿Y después de 6 horas?

\ex{3} En una pared hay dos argollas separadas 8 m entre sí. Un niño ata cada extremo de una cuerda a las argollas y se aleja de la pared hasta que la cuerda queda tensa. En ese momento, la cuerda forma ángulos de $50^\circ$ y $37^\circ$ con la pared.\begin{tasks}(2)
  \task ¿Cuánto mide la cuerda?
  \task ¿A qué distancia de la pared está el niño?
\end{tasks}

\ex{3} Uno de los ángulos de un trapecio isósceles mide $65^\circ$, los lados iguales miden 8 cm y la diagonal mide 15 cm. Determina su área.

\ex{2} Las bases de un trapecio isósceles miden 8 cm y 14 cm, y los lados iguales, 5 cm. Calcula la medida de los ángulos del trapecio.

\ex{1} ¿Qué ángulo forman entre sí las diagonales de un rectángulo de 10 cm de base y 6 cm de altura?

\ex{2} En una construcción, dos vigas de 10 m están soldadas por los extremos y forman un triángulo con otra viga de 15 m. Halla los ángulos que forman entre ellas.

\ex{2} En un romboide, los lados miden 5 cm y 8 cm y una de las diagonales mide 10 cm. Calcula la medida de sus cuatro ángulos.


% =====================================================
\seccio{9}{Fórmulas trigonométricas}

\exf{1} Escribe $15^\circ$ y $\frac{5\pi}{12}$ rad como una suma o diferencia de ángulos notables y calcula sus razones trigonométricas de manera exacta.

\ex{1} Si $\sin\alpha=0{,}6$ y $\sin\beta=0{,}8$, siendo ambos ángulos del primer cuadrante, calcula el valor de:\begin{tasks}(3)
  \task $\sin(\alpha+\beta)$
  \task $\sin(\alpha-\beta)$
  \task $\cos(\alpha+\beta)$
  \task $\cos(\alpha-\beta)$
  \task $\text{tg}\,(\alpha+\beta)$
  \task $\text{tg}\,(\alpha-\beta)$
\end{tasks}

\ex{1} Demuestra que $\sin\left(\alpha+\frac{3\pi}{2}\right)=-\cos\alpha$.

\ex{2} El seno de $20^\circ$ vale 0,34 y el coseno, 0,94. Calcula el seno y el coseno de los ángulos de $10^\circ$ y de $40^\circ$.

\ex{2} Si $\alpha$ es un ángulo del segundo cuadrante y $\sin\alpha=\frac35$, calcula las razones de $\frac{\alpha}{2}$.

\ex{3} Si $\sin\alpha=\frac35$ con $90^\circ<\alpha<180^\circ$, y $\cos\beta=-\frac14$ con $180^\circ<\beta<270^\circ$, halla: \ $\cos(\alpha+\beta)$, \ $\sin(\alpha-\beta)$, \ $\text{tg}\,2\alpha$ \ y \ $\text{tg}\,\frac{\beta}{2}$.

\ex{2} Sabiendo que $\text{tg}\,\alpha=3$, calcula las razones trigonométricas del ángulo $2\alpha$ si $\alpha$ es un ángulo:\begin{tasks}(2)
  \task Del primer cuadrante.
  \task Del tercer cuadrante.
\end{tasks}

\ex{2} Si $\cos\alpha=-\frac23$ y $90^\circ<\alpha<180^\circ$, calcula las razones trigonométricas de $\frac{\alpha}{2}$.

\ex{3} Demuestra las identidades siguientes:\begin{tasks}(2)
  \task $\dfrac{\sin\alpha-\cos\alpha}{\text{tg}\,\alpha-1}=\cos\alpha$
  \task $\text{tg}^2\alpha-\sin^2\alpha=\text{tg}^2\alpha\cdot\sin^2\alpha$
  \task $\dfrac{1+\text{cotg}\,\alpha}{\sin\alpha+\cos\alpha}=\text{cosec}\,\alpha$
  \task $\sec^2\alpha-1=\text{tg}^2\alpha$
  \task $\dfrac{1+\sin2\alpha}{1+\cos2\alpha}=\tfrac12(1+\text{tg}\,\alpha)^2$
  \task $2\,\text{tg}\,\alpha\sin^2\tfrac{\alpha}{2}+\sin\alpha=\text{tg}\,\alpha$
  \task $\cos3\alpha=4\cos^3\alpha-3\cos\alpha$
  \task $\dfrac{\text{tg}\,\alpha}{\cos2\alpha}=\text{tg}\,2\alpha-\text{tg}\,\alpha$
  \task $\text{tg}\,\alpha+\text{cotg}\,\alpha=\sec\alpha\cdot\text{cosec}\,\alpha$
  \task $\dfrac{\sin2\alpha}{1+\cos2\alpha}=\text{tg}\,\alpha$
  \task $\dfrac{1}{\text{tg}^2\alpha}=\cos^2\alpha+\left(\dfrac{\cos\alpha}{\text{tg}\,\alpha}\right)^2$
\end{tasks}

\ex{3} Demuestra las identidades siguientes:\begin{tasks}(1)
  \task $\cos\alpha\cos(\alpha-\beta)+\sin\alpha\sin(\alpha-\beta)=\cos\beta$
  \task $\sin(\alpha+\beta)\cdot\sin(\alpha-\beta)=\sin^2\alpha-\sin^2\beta$
  \task $\dfrac{\cos(\alpha+\beta)+\cos(\alpha-\beta)}{\sin(\alpha+\beta)+\sin(\alpha-\beta)}=\dfrac{1}{\text{tg}\,\alpha}$
  \task $(\cos\alpha-\cos\beta)^2+(\sin\alpha+\sin\beta)^2=4\sin^2\dfrac{\alpha+\beta}{2}$
\end{tasks}


% =====================================================
\seccio{10}{Ecuaciones trigonométricas}

\exf{1} Resuelve las ecuaciones trigonométricas siguientes. Da las soluciones de $[0^\circ,360^\circ)$ y la solución general.\begin{tasks}(3)
  \task $\sin x=\frac12$
  \task $\cos x=\frac{\sqrt3}{2}$
  \task $\text{tg}\,x=1$
  \task $\sin x=-\frac{\sqrt2}{2}$
  \task $\cos x=-\frac12$
  \task $1+\cos x=0$
  \task $\sin4x=-\frac{\sqrt3}{2}$
  \task $\cos2x=\frac{\sqrt2}{2}$
  \task $\text{tg}\,3x=-1$
  \task $\sin\frac x2=0$
  \task $\cos\frac x3=-\frac12$
  \task $\text{tg}\,\frac{3x}{4}=-\frac{\sqrt3}{3}$
\end{tasks}

\ex{2} Resuelve en $[0^\circ,360^\circ)$ las ecuaciones trigonométricas siguientes:\begin{tasks}(2)
  \task $\cos^2(2x+30^\circ)=\frac12$
  \task $\cos^2x-\sqrt3\sin x\cos x=0$
  \task $\sqrt3\cos x-\sin x=1$
  \task $\frac{1-\cos2x}{2}=2\left(\sin^2\frac x2-\cos^2\frac x2\right)$
  \task $\frac{\sin3x+\sin x}{\cos3x+\cos x}=\sqrt3$
  \task $\sin^3x-\sin x\cos^2x=0$
  \task $\sqrt3\sin x+\cos x=2$
  \task $\text{tg}^2\frac x2=1-\cos x$
  \task $\frac{\cos4x+\cos2x}{\sin4x-\sin2x}=1$
  \task $4\sin x+4\cos^2x\,\text{tg}\,x+\text{tg}\,x=0,\ \text{con}\ \text{tg}\,x\neq0$
\end{tasks}

\ex{2} Resuelve en $[0^\circ,360^\circ)$ las ecuaciones trigonométricas siguientes:\begin{tasks}(2)
  \task $3\,\text{cotg}\,x=4-\text{tg}\,x$
  \task $\sin2x\cdot\cos x=6\sin^3x$
  \task $\sin x-2\cos2x=-\frac12$
  \task $\cos(30^\circ+x)=\sin x$
  \task $\sin2x=\text{tg}\,x$
  \task $\cos3x+\cos x=0$
  \task $3\sin^2x+\cos^2x+\cos x=0$
  \task $\text{tg}\,x=\sqrt2\cos x$
  \task $\sin^2x-\cos^2x=\cos2x-7\sin x$
  \task $5\sec x-4\cos x=8$
  \task $\cos2x=1+4\sin x$
\end{tasks}


% =====================================================
\seccio{11}{Sistemas de ecuaciones trigonométricas}

\exf{2} Resuelve los sistemas de ecuaciones siguientes en el intervalo $[0^\circ,360^\circ)$:\begin{tasks}(2)
  \task $\left\{\begin{array}{l}\cos y-\sin x=1\\4\sin x\cos y+1=0\end{array}\right.$
  \task $\left\{\begin{array}{l}\cos x+\cos y=-\frac12\\\cos x\cos y=-\frac12\end{array}\right.$
  \task $\left\{\begin{array}{l}\sin x+\sin y=\frac{\sqrt3}{2}\\\sin^2x+\sin^2y=\frac34\end{array}\right.$
  \task $\left\{\begin{array}{l}x+y=\frac{\pi}{2}\\\sqrt3\cos x-\cos y=1\end{array}\right.$
  \task $\left\{\begin{array}{l}\sin x\cdot\cos y=\frac14\\\cos x\cdot\sin y=\frac14\end{array}\right.$
\end{tasks}

\ex{3} Resuelve los sistemas de ecuaciones siguientes, con $x,y\in[0,2\pi)$ en los dos primeros:\begin{tasks}(3)
  \task $\left\{\begin{array}{l}\sin x+\sin y=1\\2x+2y=\pi\end{array}\right.$
  \task $\left\{\begin{array}{l}\sin x=2\sin y\\x-y=\frac{\pi}{3}\end{array}\right.$
  \task $\left\{\begin{array}{l}4y\sin x\cos x=3\\2y\cos2x=\sqrt3\end{array}\right.$
\end{tasks}


% =====================================================
\seccio{12}{Funciones trigonométricas}

\exf{1} ¿Verdadero o falso?\begin{tasks}(1)
  \task Las funciones trigonométricas son periódicas.
  \task Las funciones $\sin$ y $\cos$ tienen un período de $2\pi$.
  \task La función $\text{tg}\,x$ tiene período $\pi$.
  \task La función $\cos x$ es como $\sin x$ desplazada $\frac{\pi}{2}$ a la izquierda.
\end{tasks}

\ex{1} ¿En qué puntos del intervalo $[0,4\pi]$ corta al eje $X$ cada una de las funciones siguientes?\begin{tasks}(3)
  \task $y=\cos\frac x2$
  \task $y=\sin(x-\pi)$
  \task $y=\cos(x+\pi)$
\end{tasks}

\ex{2} ¿Qué relación hay entre las gráficas de $y=\sin x$ e $y=\cos x$ y la de cada una de las funciones siguientes?\begin{tasks}(4)
  \task $y=\sin\left(x+\frac\pi2\right)$
  \task $y=\cos\left(x+\frac\pi2\right)$
  \task $y=\sin\left(\frac\pi2-x\right)$
  \task $y=\cos\left(\frac\pi2-x\right)$
\end{tasks}

\ex{2} Asocia a cada una de las funciones siguientes la gráfica que le corresponde:\begin{tasks}(4)
  \task $y=2\cos x$
  \task $y=\cos2x$
  \task $y=\frac12\sin x$
  \task $y=\cos\frac x2$
  \task $y=\sin2x$
  \task $y=2\sin x$
  \task $y=\frac12\cos x$
  \task $y=\sin\frac x2$
\end{tasks}

\begin{center}\begin{tikzpicture}[x=0.36cm,y=0.42cm]
\draw[muted!25,very thin] (-3.1416,-2.4) grid[xstep=1.5708,ystep=1] (12.566,2.4);
\draw[-latex,muted] (-3.4,0)--(13,0); \draw[-latex,muted] (0,-2.5)--(0,2.6);
\foreach \k/\l in {-1/-\pi,1/\pi,2/2\pi,3/3\pi,4/4\pi}{\draw[muted] (\k*3.1416,0.1)--(\k*3.1416,-0.1) node[below,font=\tiny,text=muted]{$\l$};}
\foreach \v in {-2,-1,1,2}{\draw[muted] (0.1,\v)--(-0.1,\v) node[left,font=\tiny,text=muted]{$\v$};}
\draw[navy,thick,domain=-3.1416:12.566,samples=180,smooth] plot (\x,{2*sin(\x r)});
\node[font=\small\bfseries,fill=white,inner sep=1pt] at (-2.4,2.0) {(1)};
\end{tikzpicture}\hspace{0.6cm}\begin{tikzpicture}[x=0.36cm,y=0.42cm]
\draw[muted!25,very thin] (-3.1416,-2.4) grid[xstep=1.5708,ystep=1] (12.566,2.4);
\draw[-latex,muted] (-3.4,0)--(13,0); \draw[-latex,muted] (0,-2.5)--(0,2.6);
\foreach \k/\l in {-1/-\pi,1/\pi,2/2\pi,3/3\pi,4/4\pi}{\draw[muted] (\k*3.1416,0.1)--(\k*3.1416,-0.1) node[below,font=\tiny,text=muted]{$\l$};}
\foreach \v in {-2,-1,1,2}{\draw[muted] (0.1,\v)--(-0.1,\v) node[left,font=\tiny,text=muted]{$\v$};}
\draw[navy,thick,domain=-3.1416:12.566,samples=180,smooth] plot (\x,{cos(\x/2 r)});
\node[font=\small\bfseries,fill=white,inner sep=1pt] at (-2.4,2.0) {(2)};
\end{tikzpicture}\\[4pt]
\begin{tikzpicture}[x=0.36cm,y=0.42cm]
\draw[muted!25,very thin] (-3.1416,-2.4) grid[xstep=1.5708,ystep=1] (12.566,2.4);
\draw[-latex,muted] (-3.4,0)--(13,0); \draw[-latex,muted] (0,-2.5)--(0,2.6);
\foreach \k/\l in {-1/-\pi,1/\pi,2/2\pi,3/3\pi,4/4\pi}{\draw[muted] (\k*3.1416,0.1)--(\k*3.1416,-0.1) node[below,font=\tiny,text=muted]{$\l$};}
\foreach \v in {-2,-1,1,2}{\draw[muted] (0.1,\v)--(-0.1,\v) node[left,font=\tiny,text=muted]{$\v$};}
\draw[navy,thick,domain=-3.1416:12.566,samples=180,smooth] plot (\x,{2*cos(\x r)});
\node[font=\small\bfseries,fill=white,inner sep=1pt] at (-2.4,2.0) {(3)};
\end{tikzpicture}\hspace{0.6cm}\begin{tikzpicture}[x=0.36cm,y=0.42cm]
\draw[muted!25,very thin] (-3.1416,-2.4) grid[xstep=1.5708,ystep=1] (12.566,2.4);
\draw[-latex,muted] (-3.4,0)--(13,0); \draw[-latex,muted] (0,-2.5)--(0,2.6);
\foreach \k/\l in {-1/-\pi,1/\pi,2/2\pi,3/3\pi,4/4\pi}{\draw[muted] (\k*3.1416,0.1)--(\k*3.1416,-0.1) node[below,font=\tiny,text=muted]{$\l$};}
\foreach \v in {-2,-1,1,2}{\draw[muted] (0.1,\v)--(-0.1,\v) node[left,font=\tiny,text=muted]{$\v$};}
\draw[navy,thick,domain=-3.1416:12.566,samples=180,smooth] plot (\x,{sin(\x/2 r)});
\node[font=\small\bfseries,fill=white,inner sep=1pt] at (-2.4,2.0) {(4)};
\end{tikzpicture}\\[4pt]
\begin{tikzpicture}[x=0.36cm,y=0.42cm]
\draw[muted!25,very thin] (-3.1416,-2.4) grid[xstep=1.5708,ystep=1] (12.566,2.4);
\draw[-latex,muted] (-3.4,0)--(13,0); \draw[-latex,muted] (0,-2.5)--(0,2.6);
\foreach \k/\l in {-1/-\pi,1/\pi,2/2\pi,3/3\pi,4/4\pi}{\draw[muted] (\k*3.1416,0.1)--(\k*3.1416,-0.1) node[below,font=\tiny,text=muted]{$\l$};}
\foreach \v in {-2,-1,1,2}{\draw[muted] (0.1,\v)--(-0.1,\v) node[left,font=\tiny,text=muted]{$\v$};}
\draw[navy,thick,domain=-3.1416:12.566,samples=180,smooth] plot (\x,{0.5*cos(\x r)});
\node[font=\small\bfseries,fill=white,inner sep=1pt] at (-2.4,2.0) {(5)};
\end{tikzpicture}\hspace{0.6cm}\begin{tikzpicture}[x=0.36cm,y=0.42cm]
\draw[muted!25,very thin] (-3.1416,-2.4) grid[xstep=1.5708,ystep=1] (12.566,2.4);
\draw[-latex,muted] (-3.4,0)--(13,0); \draw[-latex,muted] (0,-2.5)--(0,2.6);
\foreach \k/\l in {-1/-\pi,1/\pi,2/2\pi,3/3\pi,4/4\pi}{\draw[muted] (\k*3.1416,0.1)--(\k*3.1416,-0.1) node[below,font=\tiny,text=muted]{$\l$};}
\foreach \v in {-2,-1,1,2}{\draw[muted] (0.1,\v)--(-0.1,\v) node[left,font=\tiny,text=muted]{$\v$};}
\draw[navy,thick,domain=-3.1416:12.566,samples=180,smooth] plot (\x,{sin(2*\x r)});
\node[font=\small\bfseries,fill=white,inner sep=1pt] at (-2.4,2.0) {(6)};
\end{tikzpicture}\\[4pt]
\begin{tikzpicture}[x=0.36cm,y=0.42cm]
\draw[muted!25,very thin] (-3.1416,-2.4) grid[xstep=1.5708,ystep=1] (12.566,2.4);
\draw[-latex,muted] (-3.4,0)--(13,0); \draw[-latex,muted] (0,-2.5)--(0,2.6);
\foreach \k/\l in {-1/-\pi,1/\pi,2/2\pi,3/3\pi,4/4\pi}{\draw[muted] (\k*3.1416,0.1)--(\k*3.1416,-0.1) node[below,font=\tiny,text=muted]{$\l$};}
\foreach \v in {-2,-1,1,2}{\draw[muted] (0.1,\v)--(-0.1,\v) node[left,font=\tiny,text=muted]{$\v$};}
\draw[navy,thick,domain=-3.1416:12.566,samples=180,smooth] plot (\x,{cos(2*\x r)});
\node[font=\small\bfseries,fill=white,inner sep=1pt] at (-2.4,2.0) {(7)};
\end{tikzpicture}\hspace{0.6cm}\begin{tikzpicture}[x=0.36cm,y=0.42cm]
\draw[muted!25,very thin] (-3.1416,-2.4) grid[xstep=1.5708,ystep=1] (12.566,2.4);
\draw[-latex,muted] (-3.4,0)--(13,0); \draw[-latex,muted] (0,-2.5)--(0,2.6);
\foreach \k/\l in {-1/-\pi,1/\pi,2/2\pi,3/3\pi,4/4\pi}{\draw[muted] (\k*3.1416,0.1)--(\k*3.1416,-0.1) node[below,font=\tiny,text=muted]{$\l$};}
\foreach \v in {-2,-1,1,2}{\draw[muted] (0.1,\v)--(-0.1,\v) node[left,font=\tiny,text=muted]{$\v$};}
\draw[navy,thick,domain=-3.1416:12.566,samples=180,smooth] plot (\x,{0.5*sin(\x r)});
\node[font=\small\bfseries,fill=white,inner sep=1pt] at (-2.4,2.0) {(8)};
\end{tikzpicture}\end{center}



% =====================================================
\seccio{13}{Ampliación y problemas (ejercicios nuevos)}

\exf{1} El minutero de un reloj mide 12 cm. ¿Qué distancia recorre la punta en 25 minutos? ¿Qué área barre el minutero? (El área de un sector circular de ángulo $\alpha$ radianes es $\frac12r^2\alpha$.)

\ex{1} Calcula de manera exacta, sin calculadora: \ $\sin120^\circ\cdot\cos315^\circ+\text{tg}\,210^\circ$ \ y \ $\dfrac{\cos\frac{5\pi}{3}-\sin\frac{7\pi}{6}}{\text{tg}\,\frac{3\pi}{4}}$.

\ex{2} Simplifica las expresiones: \ $\sin(\pi-x)+\cos\left(\frac\pi2+x\right)+\sin(-x)$ \ y \ $\dfrac{\cos(\pi+x)\cdot\text{tg}\,(\pi-x)}{\sin(2\pi-x)}$.

\ex{2} Resuelve el triángulo con $a=6$ cm, $b=8$ cm y $\hat A=40^\circ$. ¿Cuántas soluciones hay?

\ex{2} Un triángulo tiene lados de 13, 14 y 15 cm. Calcula sus ángulos y su área de dos maneras: con $S=\frac12bc\sin\hat A$ y con la fórmula de Herón.

\ex{2} Dos barcos salen a la vez de un puerto con rumbos que forman un ángulo de $70^\circ$, a 18 km/h y a 24 km/h. ¿A qué distancia están al cabo de una hora y media?

\ex{2} Desde la orilla de un río vemos la copa de un árbol de la otra orilla con un ángulo de elevación de $35^\circ$. Si nos alejamos 20 m del río, perpendicularmente a la orilla, el ángulo es de $22^\circ$. Calcula la altura del árbol y la anchura del río.

\ex{2} Calcula el lado y el área de un pentágono regular inscrito en una circunferencia de 10 cm de radio.

\ex{3} Si $\text{tg}\,\alpha=2$ con $180^\circ<\alpha<270^\circ$ y $\cos\beta=-\frac35$ con $90^\circ<\beta<180^\circ$, calcula de manera exacta $\sin(\alpha+\beta)$, $\cos(\alpha-\beta)$, $\text{tg}\,(\alpha+\beta)$ y $\sin2\alpha$.

\ex{2} Demuestra que $\cos^4\alpha-\sin^4\alpha=\cos2\alpha$ y que $\dfrac{1-\cos2\alpha}{\sin2\alpha}=\text{tg}\,\alpha$.

\ex{2} Resuelve en $[0^\circ,360^\circ)$: \ $2\sin^2x+3\cos x=3$.

\ex{2} Resuelve en $[0^\circ,360^\circ)$: \ $\sin2x=\cos x$.

\ex{2} Para la función $y=2\sin\left(2x-\frac\pi2\right)+1$, indica la amplitud, el período, el recorrido y el desplazamiento respecto de $y=2\sin2x+1$. ¿En qué puntos de $[0,\pi]$ alcanza el máximo?

\vfill
\begin{center}\color{muted}\small\faCheckCircle[regular]\ \ ¡Ánimo! Comprueba tus resultados con las herramientas interactivas de los apuntes de la unidad.\end{center}

\end{document}
